\documentclass[10pt,a4paper]{article}

\usepackage[utf8]{inputenc}
\usepackage[T1]{fontenc}
\usepackage{times}
\usepackage[top=2.0cm,bottom=2.0cm,left=2.2cm,right=2.2cm]{geometry}
\usepackage{amsmath,amssymb}
\usepackage{booktabs}
\usepackage{tabularx}
\usepackage{longtable}
\usepackage{xcolor}
\usepackage{colortbl}
\usepackage{fancyhdr}
\usepackage{enumitem}
\usepackage{titlesec}
\usepackage{xurl}
\usepackage{tcolorbox}
\usepackage[hidelinks,colorlinks=true,linkcolor=blue!75!black,citecolor=teal!75!black,urlcolor=blue!75!black]{hyperref}

\linespread{1.12}
\setlength{\parindent}{0pt}
\setlength{\parskip}{4pt plus 1.5pt minus 1pt}

% Section styling
\titleformat{\section}
  {\normalfont\large\bfseries\color{blue!80!black}}{\thesection.}{0.7em}{}
\titleformat{\subsection}
  {\normalfont\normalsize\bfseries\color{teal!80!black}}{\thesubsection.}{0.7em}{}
\titleformat{\subsubsection}
  {\normalfont\small\bfseries\color{black!85}}{\thesubsubsection.}{0.7em}{}

% Header and Footer
\pagestyle{fancy}
\fancyhf{}
\fancyhead[L]{\footnotesize\color{gray}Khoirotun Nisa' (NIM: 2341720057) --- D4 Teknik Informatika Polinema}
\fancyhead[R]{\footnotesize\color{gray}Kajian Perilaku Bibit Lele \& Visi Komputer}
\fancyfoot[L]{\footnotesize\color{gray}Dosen Pembimbing: Dr. Ulla Delfana Rosiani, ST., MT.}
\fancyfoot[R]{\footnotesize\color{gray}Halaman \thepage}
\renewcommand{\headrulewidth}{0.4pt}
\renewcommand{\footrulewidth}{0.4pt}

\begin{document}

% Title Block
\begin{center}
    {\Large \textbf{\color{blue!85!black}LAPORAN KAJIAN AKADEMIK DAN SINTESIS RISET}}\\[3pt]
    {\large \textbf{\color{blue!85!black}Sistem Pemberian Pakan Adaptif Berbasis Deteksi Riak Permukaan Bibit Lele}}\\[2pt]
    {\normalsize \textbf{Analisis Kebiasaan \& Kesehatan Bibit Lele (Pakan Dibatasi vs Loss) serta Riset Mendalam Computer Vision 25 Parameter}}\\[6pt]
    \rule{\textwidth}{1pt}\\[4pt]
    \noindent
    \begin{tabularx}{\textwidth}{@{}l p{5.5cm} l X@{}}
        \textbf{Nama Mahasiswa} & : Khoirotun Nisa' & \textbf{Dosen Pembimbing} & : Dr. Ulla Delfana Rosiani, ST., MT. \\
        \textbf{NIM} & : 2341720057 & \textbf{Program Studi} & : Sarjana Terapan Teknik Informatika \\
        \textbf{Institusi} & : Politeknik Negeri Malang & \textbf{Fokus Objek} & : Bibit Ikan Lele (\textit{Clarias sp.}) \\
        \textbf{Konteks Kolam} & : Luar Ruangan (\textit{Outdoor}) + Aerator & \textbf{Arsitektur Edge} & : Raspberry Pi Camera (Top-View) \\
    \end{tabularx}\\[3pt]
    \rule{\textwidth}{1pt}
\end{center}

\vspace{-0.2cm}

\section{Konsep Sistem \& Arahan Hasil Bimbingan Dosen}
Berdasarkan arahan langsung dosen pembimbing (Dr. Ulla Delfana Rosiani, ST., MT.), fondasi sistem dirumuskan secara tegas:
\begin{enumerate}[leftmargin=*,itemsep=1pt,topsep=2pt]
    \item \textbf{Fokus Objek \& Batasan Mutlak}: Objek spesifik adalah \textbf{BIBIT IKAN LELE} (\textit{Clarias sp.}) pada fase benih/pendederan. Penelitian \textbf{TIDAK berfokus pada kualitas air} (kualitas air adalah domain kolaborator/keterbatasan lingkungan) dan \textbf{TIDAK berfokus pada diagnosis kesehatan ikan}.
    \item \textbf{Logika Sistem (\textit{Mengamati Riak $\to$ Pakan Terkontrol})}: Pemberian pakan tidak menggunakan konsep \textit{timer} kaku (``tiap jam pasti diberi pakan''). Petani memegang kendali penuh dengan menentukan kapasitas kuota harian (misal 1 kg/hari). Sistem visi komputer membagi jatah tersebut ke dalam sub-porsi (50--100 g) yang dikeluarkan mengikuti ritme riak makan bibit lele. Pakan otomatis berhenti jika kuota habis guna mencegah \textit{overfeeding}.
    \item \textbf{Tantangan Computer Vision Utama}: Kolam \textit{outdoor} menggunakan aerator terus-menerus. Gejolak aerator merupakan \textbf{BASELINE NORMAL}. Sistem tidak memakai aturan naïf ``ada gerakan = ikan lapar'', melainkan mendeteksi deviasi/perubahan temporal lonjakan riak bibit lele terhadap riak normal aerator.
    \item \textbf{Rencana Evaluasi Lapangan}: Didahului observasi baseline 1 minggu (10--15 menit tiap jam), dilanjutkan uji lapang 1 bulan membandingkan kolam sistem adaptif dengan kolam pakan manual pada parameter pertumbuhan, kelangsungan hidup (SR), dan FCR.
\end{enumerate}

\section{Output 1: Analisis Kebiasaan Bibit Lele \& Respon Pakan (Dibatasi vs Loss)}
Kajian biologis ini disintesis dari 14 literatur ilmiah fisik di folder \texttt{pakannya bibit lele} (Ikeogu et al., 2020; Tiamiyu et al., 2018; Ojelade et al., 2025; Hildebrand et al., 2025; Hossain \& Beveridge, 2001; Khormi et al., 2026; Al-salahy \& Ibrahim, 2018):

\subsection{Perilaku Makan Alami \& Respons Sensorik (\textit{Feeding Ethology})}
\textit{Clarias gariepinus} bersifat nokturnal dan sangat responsif terhadap pakan pada intensitas cahaya rendah (malam/petang) menggunakan sungut (\textit{barbels}) dan sistem gurat sisi untuk mendeteksi gelombang riak pelet (Hossain \& Beveridge, 2001). Saat terangsang pakan di permukaan, bibit lele menyambar serentak ke permukaan air, menghasilkan riak gelombang halus berdurasi singkat \textbf{3--6 frame (0,3--0,6 detik)} pada kamera 10 FPS sebelum kembali menyelam.

\subsection{Dampak Pakan Bebas / Loss (\textit{Ad Libitum / Feeding to Satiation})}
\begin{itemize}[leftmargin=*,itemsep=1pt,topsep=1pt]
    \item \textbf{Inefisiensi Konversi Pakan (FCR Buruk)}: Uji empiris \textbf{Ikeogu et al. (2020)} membuktikan bahwa pakan \textit{satiation/loss} menghasilkan FCR \textbf{$1,3 \pm 0,03$}, jauh lebih buruk dibanding pakan terkontrol (5\% biomassa) dengan FCR \textbf{$1,1 \pm 0,01$} ($p < 0,05$). Pakan \textit{loss} memboroskan 18--25\% biaya tanpa menambah laju tumbuh.
    \item \textbf{Pembusukan Pelet \& Amonia Toksik}: Pelet benih lele berkadar protein sangat tinggi (38--42\% CP). Pelet yang tidak termakan membusuk di dasar kolam, melepaskan amonia bebas ($NH_3$) dan nitrit ($NO_2^-$) yang merusak epitel insang dan memicu mortalitas benih (Tiamiyu et al., 2018; Oboh, 2022).
    \item \textbf{Distensi Lambung (\textit{Bloating})}: Pakan berlebih menyebabkan lambung bibit lele teregang melebihi kapasitas enzimatis, memicu pelet dikeluarkan kembali (\textit{regurgitation}) sebelum dicerna (Li et al., 2020).
\end{itemize}

\subsection{Dampak Pakan Dibatasi (\textit{Restricted / Rationed Feeding})}
\begin{itemize}[leftmargin=*,itemsep=1pt,topsep=1pt]
    \item \textbf{Pembatasan Terukur Optimum (5\% Biomassa/Hari)}: Ikeogu et al. (2020) dan Tiamiyu et al. (2018) menemukan bahwa takaran 5\% biomassa harian menghasilkan pertambahan bobot rata-rata ($26,11 \pm 3,65$ g) dan panjang ($6,11 \pm 0,20$ cm) tertinggi dengan kelangsungan hidup 100\%. Khormi et al. (2026) membuktikan adanya fenomena \textit{compensatory growth}, di mana efisiensi anabolisme lele meningkat pesat saat pakan diberikan tepat pada fase lapar aktif.
    \item \textbf{Bahaya Pembatasan Ekstrem / Pakan Tidak Teratur}: Studi \textbf{Ojelade et al. (2025)} membuktikan bahwa jadwal pakan acak menurunkan nilai hematokrit (PCV) dan hemoglobin lele serta memicu stres kronis dan kanibalisme antarbenih.
\end{itemize}

\textbf{Kesimpulan Integrasi Biologis}: Petani menetapkan batas atas harian (5\% biomassa) untuk mencegah pembusukan amonia dan FCR buruk, sementara sistem visi komputer menyalurkan jatah tersebut secara adaptif saat bibit lele aktif beriak di permukaan.

\section{Output 2: Riset Mendalam Computer Vision \& Ekstraksi 25 Parameter}
Berikut ekstraksi rigid 25 parameter teknis dari dua paper utama (\textbf{Hu et al., 2026} dan \textbf{Wu et al., 2024}) serta pembandingannya dengan kondisi skripsi mahasiswa:

\subsection{Tabel Komparasi 25 Parameter: Literatur Terdahulu vs Skripsi Mahasiswa}
{\scriptsize
\setlength{\tabcolsep}{2.5pt}
\begin{longtable}{@{}p{0.5cm} p{2.8cm} p{4.0cm} p{4.0cm} p{4.3cm}@{}}
\toprule
\textbf{No} & \textbf{Parameter Ekstraksi} & \textbf{Hu et al. (2026 - Nature Sci Rep)} & \textbf{Wu et al. (2024 - JMSE / Splash)} & \textbf{Kondisi Skripsi Mahasiswa} \\
\midrule
\endhead
1 & Objek Ikan & \textit{Micropterus salmoides} (Bass) & \textit{Micropterus salmoides} (Bass) & \textbf{Bibit Ikan Lele (\textit{Clarias sp.})} \\
2 & Ukuran / Fase Ikan & Dewasa (1.200 ekor teraklimatisasi) & Remaja hingga dewasa & \textbf{Benih / Pendederan (2--7 cm)} \\
3 & Lingkungan Kolam & RAS Industri Tertutup & Tangki Budidaya Terkontrol & \textbf{Kolam Luar Ruangan (\textit{Outdoor})} \\
4 & Indoor / Outdoor & Semi-indoor tertutup industri & Indoor / semi-outdoor beratap & \textbf{Outdoor Murni (Sinar Matahari/Angin)} \\
5 & Jenis Kolam & Tangki bundar $\varnothing$ 4 m, tinggi 1,5 m & Tangki persegi fiberglass & \textbf{Kolam Terpal Bundar/Kotak} \\
6 & Kondisi Permukaan Air & Bergelombang karena sirkulasi RAS & Tenang sebelum pakan disebar & \textbf{Bergolak konstan akibat aerator} \\
7 & Penggunaan Aerator & Terintegrasi hidrodinamika RAS & Dimatikan saat perekaman & \textbf{AERATOR MENYALA AKTIF KONSTAN} \\
8 & Posisi Kamera & \textit{Oblique} (miring di atas tangki) & \textit{Overhead / Top-view} (tegak lurus) & \textbf{\textit{Overhead / Top-View} (Tegak lurus)} \\
9 & Ketinggian Kamera & Sekitar 70 cm di atas air & 1,2 -- 1,5 meter di atas air & \textbf{0,8 -- 1,2 meter di atas kolam} \\
10 & Sudut Kamera & Miring $45^\circ\text{--}60^\circ$ dari permukaan & $90^\circ$ vertikal ke bawah (nadir) & \textbf{$90^\circ$ vertikal (tegak lurus kolam)} \\
11 & Resolusi Video & Resize ke $512 \times 512$ piksel & 1080p dipotong ke \textit{thumbnail} & \textbf{720p / 640$\times$480 piksel (Edge)} \\
12 & Frame Rate (FPS) & 25--30 FPS (jendela $T=5$ frame) & 30 FPS & \textbf{10 FPS (Event lele 3--6 frame)} \\
13 & Jenis Kamera & Kamera optik industri HD & Kamera digital RGB industri & \textbf{Raspberry Pi Camera Module} \\
14 & Durasi Data & Perekaman rutin multi-minggu & Sesi makan terjadwal multi-hari & \textbf{1 mgg observasi + 1 bln evaluasi} \\
15 & Waktu Pengambilan Data & Sesuai jadwal pakan otomatis & Pagi dan siang hari & \textbf{Pagi (06.00), Siang, Malam (18.00--22.00)} \\
16 & Cara Anotasi Data & 4 kelas (\textit{None, Weak, Med, Strong}) & 3 kelas (\textit{Weak, Mod, Strong Hungry}) & \textbf{2 kelas (Baseline vs Riak Lele)} \\
17 & Fitur Visual & Tekstur riak/gelombang air & \textit{Feeding splash thumbnails} & \textbf{Varians \& deviasi energi riak} \\
18 & Metode Computer Vision & CNN patch + ViT + LSTM & Deformable Attention Transformer & \textbf{Baseline Modeling + Temporal Motion} \\
19 & Model ML / Deep Learning & HVIT (ResNet18 + ViT + LSTM) & DeformAtt-ViT & \textbf{Model Temporal 1D (LSTM/Rule)} \\
20 & Input Model & Jendela 5 frame ($512\times 512\times 3$) & Potongan citra cipratan buih & \textbf{Sekuens vektor fitur gerak $T=15$} \\
21 & Output Model & 4 tingkat intensitas makan & 3 tingkat lapar + sinyal ON/OFF & \textbf{Pemicu Sub-Porsi Pakan (Trigger)} \\
22 & Pembeda Makan vs Normal & Luas riak dan kecepatan makan & Kerapatan buih cipratan air & \textbf{Lonjakan deviasi energi dari baseline} \\
23 & Penanganan Noise & Self-attention fokus area riak & Crop lokal mengabaikan pantulan & \textbf{Normalisasi Z-score baseline aerator} \\
24 & Hasil Eksperimen & Akurasi 95,10\%, Latensi 145 ms & Akurasi 95,50\%, F1 94,13\% & \textbf{Target: Akurasi $\ge 85\%$ di Raspberry Pi} \\
25 & Keterbatasan Penelitian & Model berat (28,6 M param, butuh GPU) & GAGAL PADA BENIH \& AIR AERATOR & \textbf{Tantangan komputasi edge \& cuaca} \\
\bottomrule
\end{longtable}
}

\begin{tcolorbox}[colback=blue!5!white,colframe=blue!75!black,title=\textbf{Research Gap Kunci Terhadap Wu et al. (2024) \& Hu et al. (2026)}]
Literatur penelaah mencatat bahwa metode \textit{feeding splash thumbnails} Wu et al. (2024) \textbf{tidak efektif untuk bibit ikan (\textit{fish fry})} karena massa tubuh bibit terlalu kecil untuk memercikkan air hingga menghasilkan buih putih mencolok (\textit{inconspicuous splashing}), serta mengasumsikan permukaan air awal tenang. Di sisi lain, model Hu et al. (2026) memerlukan daya komputasi GPU tinggi ($28,6$ juta parameter) dan belum diuji pada kolam terbuka dengan aerator gelembung mekanis konstan. Hal ini menjadi kebaruan (\textit{novelty}) skripsi ini: \textbf{mendeteksi riak gelombang halus bibit lele di atas baseline turbulensi aerator pada kolam outdoor menggunakan sistem edge efisien}.
\end{tcolorbox}

\subsection{Karakteristik Metode Computer Vision untuk Kolam Ber-aerator}
\begin{enumerate}[leftmargin=*,itemsep=1pt]
    \item \textbf{Frame Difference / Temporal Difference}: Sangat ringan (>30 FPS di Raspberry Pi), sangat tepat sebagai penyaring awal (\textit{first-stage filter}), namun rentan terhadap fluktuasi intensitas sinar matahari.
    \item \textbf{Dense Optical Flow (Farnebäck)}: Menghasilkan vektor magnitudo kecepatan riak, tetapi komputasinya sangat lambat di CPU Raspberry Pi 4 (<2 FPS) dan mengasumsikan permukaan air diam.
    \item \textbf{Motion Intensity Index}: Menghitung magnitudo energi kinetik agregat seluruh ROI kolam. Sangat efisien, independen dari bentuk ikan, dan ideal dijadikan sinyal sekuensial 1D.
    \item \textbf{Background Subtraction (MOG2 / KNN)}: Memisahkan latar depan dinamis, namun gelembung aerator yang bergerak konstan berisiko selalu masuk sebagai \textit{foreground} jika parameter \textit{learning rate} tidak diadaptasi.
    \item \textbf{Ripple / Splash Detection}: Cipratan bibit lele tidak menghasilkan buih putih seperti ikan dewasa Wu et al., sehingga deteksi harus difokuskan pada frekuensi dan kontras riak gelombang air.
    \item \textbf{Model Temporal (LSTM / GRU / 1D-CNN)}: Sangat unggul untuk membedakan sifat periodik aerator dengan sifat impulsif riak bibit lele.
\end{enumerate}

\subsection{Spesifikasi Rancang Bangun Model Temporal (LSTM)}
Jika LSTM digunakan pada skripsi ini, arsitekturnya dirancang sebagai berikut:
\begin{itemize}[leftmargin=*,itemsep=1.5pt]
    \item \textbf{Input Sequence}: Bukan raw video piksel berat, melainkan \textbf{vektor fitur gerak 1D terkompresi} berdimensi $(B, T, D)$ dengan $B=1$, $T=15$ frame ($1,5$ detik pada 10 FPS), dan $D=4$ fitur:
    \begin{enumerate}[label=(\alph*),itemsep=0pt]
        \item $M_t$: Indeks intensitas gerak rata-rata (\textit{motion intensity}) dari selisih frame terbobot.
        \item $V_t$: Varians spasial sebaran gerak (membedakan riak terpusat di titik aerator vs menyebar di kolam).
        \item $A_t$: Luas persentase piksel aktif di atas ambang batas dinamis.
        \item $\Delta M_t$: Gradien laju perubahan energi antar-frame ($M_t - M_{t-1}$).
    \end{enumerate}
    \item \textbf{Pembentukan Sequence}: \textit{Sliding window} panjang $T=15$ frame dengan \textit{stride} 3 frame (overlap 70\%). Rentang $1,5$ detik mencakup riwayat sebelum event (\textit{pre-event baseline}), puncak riak bibit lele (3--6 frame), dan fase relaksasi setelahnya.
    \item \textbf{Target Klasifikasi Biner}: Label \texttt{0} (Baseline Aerator Normal) dan Label \texttt{1} (Lonjakan Riak Bibit Lele).
    \item \textbf{Mengapa Model Temporal Diperlukan?}: Gejolak aerator bersifat \textbf{periodik/siklis konstan}, sedangkan riak bibit lele bersifat \textbf{impulsif dan acak}. Model spasial murni (frame tunggal) mudah tertipu oleh letupan buih aerator sesaat (\textit{false alarm}). LSTM memanfaatkan memori sekuensial untuk memastikan adanya deviasi tren energi kinetik nyata.
    \item \textbf{Penanganan Baseline Aerator}: Sistem menghitung rata-rata bergerak ($\mu_{\text{base}}$) dan deviasi standar ($\sigma_{\text{base}}$) riak aerator, lalu menormalkan fitur melalui $Z_t = \frac{M_t - \mu_{\text{base}}}{\sigma_{\text{base}}}$. Ditambahkan periode jeda (\textit{refractory cooldown} 3--5 menit) pasca-pakan agar riak sisa makan tidak memicu pakan berulang.
\end{itemize}

\section{Rancangan Evaluasi Lapangan (1 Bulan vs Kolam Manual)}
\begin{enumerate}[leftmargin=*,itemsep=1pt]
    \item \textbf{Tahap Observasi Awal (1 Minggu)}: Perekaman video 10--15 menit tiap jam untuk membangun profil variabilitas baseline aerator dan ritme alami riak bibit lele pada kondisi pencahayaan luar ruangan (pagi, siang, petang, malam).
    \item \textbf{Tahap Pengujian Lapangan (1 Bulan)}:
    \begin{itemize}[itemsep=0.5pt]
        \item \textbf{Kolam Perlakuan (Sistem Adaptif)}: Pakan diberikan berbasis deteksi riak kamera + pembatasan kuota harian petani.
        \item \textbf{Kolam Kontrol (Manual)}: Pakan diberikan manual sesuai jadwal konvensional petani tanpa alat.
    \end{itemize}
    \item \textbf{Parameter Evaluasi Lapangan Terukur}:
    \begin{itemize}[itemsep=0.5pt]
        \item Pertumbuhan bobot mutlak (\textit{Mean Weight Gain}) dan laju pertumbuhan spesifik (\textit{Specific Growth Rate} / SGR).
        \item Pertambahan panjang tubuh bibit lele (\textit{Length Gain}).
        \item Tingkat kelangsungan hidup benih (\textit{Survival Rate} / SR).
        \item Total kuota pakan yang dikonsumsi (kg) dan Rasio Konversi Pakan (\textit{Feed Conversion Ratio} / FCR).
        \item Efisiensi waktu dan beban kerja fisik petani.
    \end{itemize}
\end{enumerate}

\vspace{0.2cm}
\noindent\rule{\textwidth}{0.5pt}
\begin{center}
    \footnotesize\color{gray}
    Dokumen ini merupakan sintesis resmi persiapan skripsi Jurusan Teknologi Informasi Politeknik Negeri Malang (2026).\\
    Diselaraskan berdasarkan catatan konsultasi dengan Dr. Ulla Delfana Rosiani, ST., MT.
\end{center}

\end{document}
